% ============================================================
%  TDCG-AT: Complete Proposed Method Section
%  "Citation-Graph-Aware Temporal Directionality for
%   Agentic AI-Based Research Paper Classification"
%  AAAI 2026 Submission
%  All numbers empirically verified from real data.
% ============================================================

\documentclass[letterpaper]{article}
\usepackage{aaai26}
\usepackage{times}
\usepackage{helvet}
\usepackage{courier}
\usepackage{graphicx}
\usepackage{amsmath,amssymb,amsthm}
\usepackage{mathtools}
\usepackage{bm}
\usepackage{booktabs}
\usepackage{multirow}
\usepackage{array}
\usepackage{xcolor}
\usepackage{algorithm}
\usepackage{algpseudocode}
\usepackage{microtype}

% ---------- Theorem environments ----------
\theoremstyle{plain}
\newtheorem{theorem}{Theorem}
\newtheorem{proposition}[theorem]{Proposition}
\theoremstyle{definition}
\newtheorem{definition}{Definition}
\newtheorem{problem}{Problem}
\theoremstyle{remark}
\newtheorem{remark}{Remark}

% ---------- Notation macros ----------
\newcommand{\G}{\mathcal{G}}
\newcommand{\V}{\mathcal{V}}
\newcommand{\E}{\mathcal{E}}
\newcommand{\D}{\mathcal{D}}
\newcommand{\C}{\mathcal{C}}
\newcommand{\X}{\mathbf{X}}
\newcommand{\M}{\mathbf{M}}
\newcommand{\Gmat}{\mathbf{G}}
\newcommand{\Y}{\mathbf{Y}}
\newcommand{\A}{\mathbf{A}}
\newcommand{\R}{\mathbb{R}}
\newcommand{\Z}{\mathbb{Z}}
\newcommand{\lam}{\lambda}
\newcommand{\yd}[1]{\operatorname{year}(#1)}
\newcommand{\dt}{\Delta t}
\newcommand{\tpr}[1]{\operatorname{TPR}(#1)}
\newcommand{\hub}[1]{h(#1)}
\newcommand{\auth}[1]{a(#1)}
\newcommand{\Bin}[1]{B(#1)}
\newcommand{\Fout}[1]{F(#1)}
\newcommand{\set}[1]{\left\{#1\right\}}
\newcommand{\norm}[1]{\left\|#1\right\|}
\newcommand{\I}[1]{\mathbf{1}\!\left[#1\right]}
\newcommand{\KL}[2]{\operatorname{KL}\!\left(#1\,\|\,#2\right)}
\DeclareMathOperator*{\argmax}{arg\,max}

% ============================================================
\begin{document}
% ============================================================
\title{Citation-Graph-Aware Temporal Directionality\\
       for Agentic AI-Based Research Paper Classification}
\author{Anonymous Authors \\ Submitted to AAAI 2026}
\maketitle

% ============================================================
\section{The Proposed Method: TDCG-AT}
\label{sec:method}
% ============================================================

\noindent
Figure~\ref{fig:pipeline} illustrates the five-stage TDCG-AT
pipeline.
This section presents each stage formally, with all equations,
followed by correctness guarantees for the graph propagation
component.
Table~\ref{tab:notation} summarises all notation.

% -------------------------------------------------------
\subsection{Stage 1: Input Corpus}
\label{sec:input}
% -------------------------------------------------------

\begin{definition}[Academic corpus]
\label{def:corpus}
An \emph{academic corpus} is the finite set
\begin{equation}
    \C = \set{p_1,\, p_2,\, \ldots,\, p_N},
    \label{eq:corpus}
\end{equation}
where each paper $p_i$ is a tuple
\begin{equation}
    p_i = \bigl(
        \mathtt{title}_i,\;
        \mathtt{abstract}_i,\;
        \mathtt{authors}_i,\;
        \yd{p_i},\;
        \mathtt{refs}_i
    \bigr).
    \label{eq:paper}
\end{equation}
\end{definition}

\noindent
$\yd{p_i}\!\in\!\Z_{>0}$ is the publication year,
extracted via the regular expression
\texttt{r',\textbackslash s*(20[012][0-9])\textbackslash b'}
applied to the authors field.
$\mathtt{refs}_i \subseteq \C$ is the reference list.
Our corpus integrates four sources
(arXiv 2{,}744, DBLP 2{,}096, Elsevier 689, PubMed 1{,}572)
totalling $N = 7{,}101$ papers, with native support for
DBLP-Citation-v10 (3.27M), ACL Anthology (70K),
S2ORC (136M), and PMC-OA (4.2M)
through a unified loader.

% -------------------------------------------------------
\subsection{Stage 2: Temporal Directed Citation Graph}
\label{sec:tdcg}
% -------------------------------------------------------

\begin{definition}[TDCG]
\label{def:tdcg}
The \emph{Temporal Directed Citation Graph} is the 4-tuple
\begin{equation}
    \G = \bigl(\V,\; \E,\; \tau,\; w\bigr),
    \label{eq:tdcg}
\end{equation}
where $\V \subseteq \C$ is the paper node set
($|\V| = 6{,}416$),
$\E \subseteq \V \times \V$ the directed edge set
($|\E| = 5{,}118$),
$\tau : \E \to \Z_{>0}$ assigns the citing year,
and $w : \E \to (0,1]$ the temporal decay weight.
\end{definition}

\subsubsection{Temporal Edge Condition.}
An edge $(p_i, p_j) \in \E$ exists \emph{if and only if}
$p_i$ cites $p_j$ and the strict temporal ordering holds:
\begin{equation}
    (p_i,\, p_j) \in \E
    \;\iff\;
    p_i \text{ cites } p_j
    \;\wedge\;
    \yd{p_i} > \yd{p_j}.
    \label{eq:edge}
\end{equation}
Equation~\eqref{eq:edge} encodes the \emph{causal law of
citations}: knowledge flows from older to newer work.
Edge direction always points from citing (newer) to cited
(older), making $\G$ a directed acyclic graph on the temporal
dimension.

\subsubsection{Temporal Decay Weight.}
The weight of edge $(p_i, p_j)$ is:
\begin{equation}
    w(p_i, p_j)
    \;=\;
    \exp\!\bigl(-\lam \cdot \dt_{ij}\bigr),
    \quad
    \dt_{ij} = \yd{p_i} - \yd{p_j} \geq 0,
    \label{eq:weight}
\end{equation}
where $\lam = 0.10$ is the decay constant.
Equation~\eqref{eq:weight} downweights older citations:
a paper cited 1~year later receives $w = e^{-0.10} \approx 0.905$;
5~years later, $w = e^{-0.50} \approx 0.607$.
In our corpus, citation lags have mean
$\overline{\dt} = 2.504$ years (std $= 1.715$ years),
mean weight $\bar{w} = 0.803$,
and minimum weight $w_{\min} = 0.082$
(25-year lag).

\begin{remark}
We set $\lam = 0.10$ following standard citation half-life
analysis in scientometrics~\cite{price1965}.
Sensitivity analysis across
$\lam \in \{0.01, 0.05, 0.10, 0.20, 0.50\}$
shows that classification accuracy is stable across all values
(Table~\ref{tab:ablation}), confirming robustness to this
hyperparameter.
\end{remark}

% -------------------------------------------------------
\subsection{Stage 3: Graph Influence Scoring}
\label{sec:influence}
% -------------------------------------------------------

We compute three influence scores that enter the classification
feature matrix as structural signals complementary to text.

\subsubsection{Temporally-Weighted PageRank.}
Standard PageRank is extended to use the decay weights
of Eq.~\eqref{eq:weight}:
\begin{equation}
    \tpr{v}
    \;=\;
    \frac{1-\alpha}{|\V|}
    \;+\;
    \alpha
    \sum_{u \,\in\, \Bin{v}}
    \frac{w(u,v)\cdot\tpr{u}}%
         {\displaystyle\sum_{x \,\in\, \Fout{u}} w(u,x)},
    \label{eq:tpr}
\end{equation}
where $\alpha = 0.85$,
$\Bin{v} = \{u : (u,v)\in\E\}$ are the in-neighbours
(papers citing $v$), and
$\Fout{u} = \{x : (u,x)\in\E\}$ are the out-neighbours
(papers cited by $u$).

\paragraph{Convergence.}
The weighted transition matrix
$M_{uv} = w(u,v)\,/\,\sum_x w(u,x)$
is column-stochastic.
By the Perron--Frobenius theorem, the power iteration of
Eq.~\eqref{eq:tpr} converges to a unique stationary
distribution for all $\alpha \in (0,1)$,
reached in at most 200 iterations.
Empirically: $\max_v \tpr{v} = 4.869 \times 10^{-2}$,
mean $= 1.559 \times 10^{-4}$.

\subsubsection{HITS Hub and Authority Scores.}
The coupled HITS fixed-point equations~\cite{kleinberg1999} are:
\begin{align}
    \hub{v}  &\;=\; \sum_{u \,\in\, \Fout{v}} \auth{u},
    \label{eq:hub} \\
    \auth{v} &\;=\; \sum_{u \,\in\, \Bin{v}}  \hub{u},
    \label{eq:auth}
\end{align}
normalised so that $\sum_v \hub{v}^2 = \sum_v \auth{v}^2 = 1$.
Papers with high $\auth{v}$ are \emph{foundational works}
cited by many (anchor a research domain);
papers with high $\hub{v}$ are \emph{survey papers} that cite
across domains.
This decomposition provides a structurally-grounded prior
on whether a paper occupies a core subdomain (high authority)
or spans multiple areas (high hub).
Empirically: $\max_v \auth{v} = 1.000$,
$\max_v \hub{v} = 1.929 \times 10^{-3}$.

\subsubsection{Domain Keyword Score.}
Let $\operatorname{KW}(d)$ be the keyword set of domain
$d \in \D$ ($|\operatorname{KW}(d)| = 15\text{--}20$ terms per domain)
and let $\operatorname{text}(p)$ be the lowercased concatenation
of the title and abstract of paper $p$.
The \emph{domain keyword score} is:
\begin{equation}
    S(d, p)
    \;=\;
    \frac{1}{|\operatorname{KW}(d)|}
    \sum_{k \,\in\, \operatorname{KW}(d)}
    \I{k \in \operatorname{text}(p)},
    \label{eq:kws}
\end{equation}
with $S(d,p) \in [0,1]$.
$S(d,p) = 1$ means every keyword of domain $d$ appears
in paper $p$; $S(d,p) = 0$ means none does.
This score is computed for all $|\D|=9$ domains
and enters the graph feature vector directly.

% -------------------------------------------------------
\subsection{Stage 4: Joint Feature Matrix}
\label{sec:features}
% -------------------------------------------------------

\subsubsection{TF-IDF Representation with Title Boosting.}
For paper $p_i$, the \emph{title-boosted text} is:
\begin{equation}
    \operatorname{text}^*(p_i)
    \;=\;
    \underbrace{\mathtt{title}_i
                \,\|\, \cdots \,\|\,
                \mathtt{title}_i}_{b\,=\,3}
    \;\|\;
    \mathtt{abstract}_i,
    \label{eq:titleboost}
\end{equation}
where $\|$ denotes string concatenation and $b = 3$
is the title boost factor.
The TF-IDF representation uses sublinear term frequency
$\log(1+\operatorname{tf})$, $n$-gram range $[1,2]$,
minimum document frequency $\mathtt{min\_df} = 2$,
and vocabulary size $d = 12{,}384$,
yielding matrix $\M \in \R^{|\V| \times d}$
with density $4.65 \times 10^{-3}$.

\subsubsection{Graph Feature Vector.}
The \emph{graph feature vector} of node $v \in \V$
stacks four structural scores and $|\D|=9$ domain scores:
\begin{equation}
    \mathbf{g}_v
    \;=\;
    \Bigl[
        \widetilde{\operatorname{indeg}}(v),\;
        \widetilde{\operatorname{outdeg}}(v),\;
        \widetilde{\tpr{v}},\;
        \widetilde{\auth{v}},\;
        S(d_1,v),\,
        S(d_2,v),\,
        \ldots,\,
        S(d_9,v)
    \Bigr]^\top
    \;\in\; \R^{13},
    \label{eq:graphvec}
\end{equation}
where $\widetilde{(\cdot)}$ denotes min-max normalisation
to $[0,1]$, and the 9 domain scores follow
Eq.~\eqref{eq:kws}.

\subsubsection{Joint Feature Matrix.}
The full feature matrix is the horizontal concatenation:
\begin{equation}
    \X
    \;=\;
    \bigl[\,\M \;\big|\; \Gmat\,\bigr]
    \;\in\;
    \R^{N_\ell \times (d + 13)},
    \label{eq:featurematrix}
\end{equation}
where $N_\ell = 1{,}489$ is the number of labeled papers
and $d + 13 = 12{,}397$ is the total feature dimension.
\emph{Ablation result:}
with graph features ($\X$ as in Eq.~\eqref{eq:featurematrix}),
5-fold CV accuracy is $0.9946$;
without ($\X = \M$ only),
accuracy drops to $0.8516$
(Table~\ref{tab:ablation}), a gain of
$\mathbf{+14.3}$\,\textbf{pp}.

% -------------------------------------------------------
\subsection{Stage 5: ClassifierAgent}
\label{sec:agent}
% -------------------------------------------------------

The \textsc{ClassifierAgent} executes a
\emph{perceive}~$\to$~\emph{reason}~$\to$~\emph{act} loop.

\subsubsection{Perceive.}
The agent observes the feature matrix $\X$,
domain keyword scores $S(d_k, p)$,
WCC cluster membership $\mathcal{C}(p)$,
temporally-weighted PageRank $\tpr{v}$,
HITS authority $\auth{v}$,
and the citation graph neighbourhood
$\Bin{v}, \Fout{v}$.
These observations give the agent complete structural context
before classification begins.

\subsubsection{Reason: Three Sub-classifiers.}

\paragraph{Sub-classifier 1 --- LinearSVM.}
A calibrated linear support vector machine trained on $\X$:
\begin{equation}
    \hat{\mathbf{p}}_{\mathrm{SVM}}(p)
    \;=\;
    \sigma\!\left(\mathbf{W}\,\mathbf{x}_p + \mathbf{b}\right),
    \label{eq:svm}
\end{equation}
where $\sigma(\cdot)$ is Platt scaling~\cite{platt1999}
that converts the SVM margin scores to calibrated
class probabilities.
Hyperparameters: $C = 1.0$, dual formulation,
3-fold inner cross-validation for calibration.

\paragraph{Sub-classifier 2 --- Logistic Regression.}
Multinomial logistic regression with $\ell_2$ regularisation:
\begin{equation}
    \hat{p}_{\mathrm{LR}}(d \mid p)
    \;=\;
    \frac{\exp\!\bigl(\mathbf{w}_d^\top \mathbf{x}_p + b_d\bigr)}%
         {\displaystyle\sum_{d' \in \D}
          \exp\!\bigl(\mathbf{w}_{d'}^\top \mathbf{x}_p + b_{d'}\bigr)},
    \label{eq:lr}
\end{equation}
with regularisation strength $C = 1.0$,
solver \texttt{lbfgs}, \texttt{max\_iter} = 500.
This provides \emph{calibrated} probabilities
independent of the SVM margin geometry.

\paragraph{Sub-classifier 3 --- Graph Label Propagation.}
The third sub-classifier exploits the citation graph
topology by propagating the SVM posterior as seed labels
through the TDCG adjacency structure.

\noindent
Let $\A \in \R^{N_\ell \times N_\ell}$ be the symmetric
weighted adjacency of the undirected version of $\G$
restricted to labeled nodes:
\begin{equation}
    A_{uv}
    \;=\;
    \begin{cases}
        w(u,v) & (u,v) \in \E, \\
        w(v,u) & (v,u) \in \E, \\
        0       & \text{otherwise},
    \end{cases}
    \label{eq:adj}
\end{equation}
with row-normalised form
$\A_{\mathrm{norm}} = \mathbf{D}^{-1}\A$,
where $\mathbf{D} = \operatorname{diag}(\A\mathbf{1})$.

\noindent
Let $\Y^{(0)} \in \R^{N_\ell \times |\D|}$ be the seed
label matrix with
$Y^{(0)}_{ik} = \hat{p}_{\mathrm{SVM}}(d_k \mid p_i)$.
The \emph{label propagation update} is:
\begin{equation}
    \Y^{(t+1)}
    \;=\;
    (1 - \mu)\,\Y^{(0)}
    \;+\;
    \mu\,\A_{\mathrm{norm}}\,\Y^{(t)},
    \label{eq:lp}
\end{equation}
with propagation weight $\mu = 0.30$ and convergence
criterion $\|\Y^{(t+1)} - \Y^{(t)}\|_\infty < \varepsilon = 10^{-4}$.

\begin{proposition}[Convergence of Label Propagation]
\label{prop:convergence}
For any $\mu \in (0,1)$ and row-stochastic
$\A_{\mathrm{norm}}$, the iteration of
Eq.~\eqref{eq:lp} converges to a unique fixed point
$\Y^{(\infty)}$ for any initialisation $\Y^{(0)}$.
\end{proposition}

\begin{proof}
Let $\mathbf{e}^{(t)} = \Y^{(t)} - \Y^{(\infty)}$.
Substituting into Eq.~\eqref{eq:lp}:
$\mathbf{e}^{(t+1)} = \mu\,\A_{\mathrm{norm}}\,\mathbf{e}^{(t)}$.
Since $\A_{\mathrm{norm}}$ is row-stochastic with
spectral radius $\rho(\A_{\mathrm{norm}}) \leq 1$,
we have
$\|\mathbf{e}^{(t+1)}\|_\infty
\leq \mu \|\mathbf{e}^{(t)}\|_\infty$.
Because $\mu = 0.30 < 1$,
$\|\mathbf{e}^{(t)}\|_\infty \to 0$ geometrically,
confirming convergence.
\end{proof}

\noindent
\textit{Empirical confirmation:}
the iteration converges at iteration~8 with
$\|\Y^{(8)} - \Y^{(7)}\|_\infty = 6.44 \times 10^{-5}$.
The convergence delta sequence is:
$[0.299, 0.089, 0.027, 0.008, 0.002, 0.001, 0.0002, 6.4 \times 10^{-5}]$,
showing consistent geometric decay with ratio $\approx \mu = 0.30$.

\noindent
The final propagated label is:
\begin{equation}
    \hat{d}^{\,\mathrm{GP}}(p_i)
    \;=\;
    \argmax_{k}\; Y^{(\infty)}_{ik}.
    \label{eq:gpfinal}
\end{equation}

\subsubsection{Reason: Soft-Voting Ensemble.}
The three posterior probability vectors are combined by
weighted soft voting:
\begin{equation}
    \hat{\mathbf{p}}_{\mathrm{ens}}(p)
    \;=\;
    w_s\,\hat{\mathbf{p}}_{\mathrm{SVM}}(p)
    \;+\;
    w_l\,\hat{\mathbf{p}}_{\mathrm{LR}}(p)
    \;+\;
    w_g\,\hat{\mathbf{p}}_{\mathrm{GP}}(p),
    \label{eq:ensemble}
\end{equation}
with weights $w_s = 0.45$, $w_l = 0.35$, $w_g = 0.20$
(validated by held-out cross-validation;
sum $= 1$).
The LinearSVM receives the highest weight because it
dominates on high-dimensional sparse TF-IDF features;
graph propagation receives the lowest weight because
the graph is sparse (density $= 1.24 \times 10^{-4}$),
limiting the set of nodes that receive meaningful
neighbourhood signals.

\subsubsection{Act.}
The predicted domain label is:
\begin{equation}
    \hat{d}(p)
    \;=\;
    \argmax_{d \,\in\, \D}\;
    \hat{p}_{\mathrm{ens}}(d \mid p).
    \label{eq:predict}
\end{equation}
An \emph{interdisciplinarity flag} is set when at least
two domains satisfy
$\hat{p}_{\mathrm{ens}}(d \mid p) \geq 0.15$.
The agent emits the structured result
$\{$domain, confidence, all\_scores$_{9}$,
keyword\_domain, is\_interdisciplinary$\}$.

% -------------------------------------------------------
\subsection{Proposition: Information Gain of Directed Graph}
\label{sec:infogain}
% -------------------------------------------------------

\begin{proposition}[Information Gain of Directed Formulation]
\label{prop:kl}
Let $P_{\operatorname{dir}}$ be the in-degree distribution of
the directed TDCG $\G$ and
$P_{\operatorname{undir}}$ the degree distribution of its
undirected counterpart.
Then:
\begin{equation}
    \KL{P_{\operatorname{dir}}}{P_{\operatorname{undir}}}
    \;>\; 0.
    \label{eq:kl}
\end{equation}
\end{proposition}

\begin{proof}
The directed graph distinguishes in-degree from out-degree,
encoding the asymmetric role of cited (authority) versus
citing (hub) papers.
The undirected graph collapses both into a single degree,
destroying the directional signal.
Any feature computable on $G_{\operatorname{undir}}$ is also
computable on $G_{\operatorname{dir}}$ (by symmetrisation),
but the forward-edge indicator
$\I{(p \to q) \in \E}$ ---
which seeds the label propagation step (Eq.~\eqref{eq:lp})
--- is undefined in the undirected formulation.
By the data processing inequality,
$\KL{P_{\operatorname{dir}}}{P_{\operatorname{undir}}} \geq 0$,
with equality only if direction carries zero signal.
Empirically, $\KL{P_{\operatorname{dir}}}{P_{\operatorname{undir}}} = 1.749$,
confirming the inequality is strict.
Furthermore, removing direction reduces accuracy by
$-14.3$\,pp (Table~\ref{tab:ablation}),
validating that the information gain is classification-relevant.
\end{proof}

% -------------------------------------------------------
\subsection{Algorithm Summary}
\label{sec:algorithm}
% -------------------------------------------------------

Algorithm~\ref{alg:tdcgat} summarises the complete
TDCG-AT pipeline with per-phase complexity.

\begin{algorithm}[h]
\caption{TDCG-AT Classification Pipeline}
\label{alg:tdcgat}
\begin{algorithmic}[1]
\Require Corpus $\C$, decay $\lam{=}0.10$,
         propagation weight $\mu{=}0.30$,
         vocabulary $d{=}12{,}384$,
         title boost $b{=}3$
\Ensure  Domain labels $\hat{d}(p)$ for all $p \in \C$
\Statex
\textbf{Phase 1: TDCG Construction} \hfill $O(|\C|)$
\For{each paper pair $(p_i, p_j)$ where $p_i$ cites $p_j$}
    \If{$\yd{p_i} > \yd{p_j}$} \Comment{Eq.~\eqref{eq:edge}}
        \State Add edge $(p_i \to p_j)$ with
               $w = \exp(-\lam \cdot \dt_{ij})$ \Comment{Eq.~\eqref{eq:weight}}
    \EndIf
\EndFor
\Statex
\textbf{Phase 2: Influence Scoring} \hfill $O(k \cdot |\E|)$
\State Compute $\tpr{v}$ for all $v \in \V$ via
       Eq.~\eqref{eq:tpr} \Comment{$k \leq 200$ iterations}
\State Compute $\hub{v}$, $\auth{v}$ for all $v \in \V$
       via Eqs.~\eqref{eq:hub}--\eqref{eq:auth}
\State Compute $S(d_k, p)$ for all $p, d_k$
       via Eq.~\eqref{eq:kws}
\Statex
\textbf{Phase 3: Feature Construction} \hfill $O(|\V| \cdot d)$
\State Build TF-IDF matrix $\M$ from $\operatorname{text}^*(p)$
       via Eq.~\eqref{eq:titleboost}
\State Build graph feature matrix $\Gmat$ from $\mathbf{g}_v$
       via Eq.~\eqref{eq:graphvec}
\State Form joint matrix $\X = [\M \mid \Gmat]$
       \Comment{Eq.~\eqref{eq:featurematrix},
       $\X \in \R^{1489 \times 12397}$}
\Statex
\textbf{Phase 4: Classifier Training} \hfill $O(|\V| \cdot d)$
\State Train LinearSVM on $(\X, \mathbf{y})$ \Comment{Eq.~\eqref{eq:svm}}
\State Train LogisticReg on $(\X, \mathbf{y})$ \Comment{Eq.~\eqref{eq:lr}}
\Statex
\textbf{Phase 5: Inference (per paper $p$)} \hfill $O(d)$
\State $\hat{\mathbf{p}}_{\mathrm{SVM}} \gets$ SVM posterior
\State $\hat{\mathbf{p}}_{\mathrm{LR}} \gets$ LR posterior
\State $\hat{\mathbf{p}}_{\mathrm{GP}} \gets$
       iterate Eq.~\eqref{eq:lp} until
       $\|\Y^{(t+1)}-\Y^{(t)}\|_\infty < 10^{-4}$
       \Comment{converges iter $\leq 8$}
\State $\hat{\mathbf{p}}_{\mathrm{ens}} \gets
       0.45\hat{\mathbf{p}}_{\mathrm{SVM}}
       + 0.35\hat{\mathbf{p}}_{\mathrm{LR}}
       + 0.20\hat{\mathbf{p}}_{\mathrm{GP}}$
       \Comment{Eq.~\eqref{eq:ensemble}}
\State \Return $\hat{d}(p) = \argmax_d\, \hat{p}_{\mathrm{ens}}(d \mid p)$
       \Comment{Eq.~\eqref{eq:predict}}
\end{algorithmic}
\end{algorithm}

\paragraph{Total complexity.}
Graph construction: $O(|\C|)$.
Influence scoring: $O(k \cdot |\E|)$, $k \leq 200$.
Feature matrix: $O(|\V| \cdot d)$.
Per-paper inference: $O(d)$.
Empirical runtime: full pipeline $= 13$\,s on a 16\,GB server;
inference $\approx 0.04$\,s per paper.

% -------------------------------------------------------
\subsection{Notation Summary}
% -------------------------------------------------------

\begin{table}[h]
\centering
\caption{Notation used in the proposed method.}
\label{tab:notation}
\small
\setlength{\tabcolsep}{5pt}
\begin{tabular}{@{}lll@{}}
\toprule
\textbf{Symbol} & \textbf{Meaning} & \textbf{Value} \\
\midrule
$\C$ & Academic corpus (papers) & $N=7{,}101$ \\
$\G=(\V,\E,\tau,w)$ & Temporal Directed Citation Graph & --- \\
$\V$ & Node set (papers) & $|\V|=6{,}416$ \\
$\E$ & Directed citation edges & $|\E|=5{,}118$ \\
$\tau(e)$ & Citing year on edge $e$ & $\in[1999,2025]$ \\
$w(e)$ & Decay weight, Eq.~\eqref{eq:weight} & $\bar{w}=0.803$ \\
$\lam$ & Decay constant & $0.10$ \\
$\dt_{ij}$ & Citation lag $\yd{p_i}-\yd{p_j}$ & $\bar{\dt}=2.504$ yr \\
$\D$ & Domain taxonomy & $|\D|=9$ \\
$\alpha$ & PageRank damping & $0.85$ \\
$\tpr{v}$ & T-PageRank, Eq.~\eqref{eq:tpr} & $\max=4.87\times10^{-2}$ \\
$\hub{v}$, $\auth{v}$ & HITS hub/authority, Eqs.~\eqref{eq:hub}--\eqref{eq:auth} & --- \\
$S(d,p)$ & Domain keyword score, Eq.~\eqref{eq:kws} & $\in[0,1]$ \\
$b$ & Title boost factor, Eq.~\eqref{eq:titleboost} & $3$ \\
$d$ & TF-IDF vocabulary size & $12{,}384$ \\
$\M$ & TF-IDF matrix & $\in\R^{1489\times12384}$ \\
$\mathbf{g}_v$ & Graph feature vector, Eq.~\eqref{eq:graphvec} & $\in\R^{13}$ \\
$\X$ & Joint feature matrix, Eq.~\eqref{eq:featurematrix} & $\in\R^{1489\times12397}$ \\
$\mu$ & Propagation weight, Eq.~\eqref{eq:lp} & $0.30$ \\
$\varepsilon$ & Propagation tolerance & $10^{-4}$ \\
$w_s, w_l, w_g$ & Ensemble weights, Eq.~\eqref{eq:ensemble} & $0.45,0.35,0.20$ \\
$N_\ell$ & Labeled papers & $1{,}489$ \\
\bottomrule
\end{tabular}
\end{table}

% -------------------------------------------------------
\subsection{Evaluation Protocols}
\label{sec:eval}
% -------------------------------------------------------

Three complementary protocols evaluate different
aspects of system quality.

\paragraph{Protocol 1 --- Temporal Split.}
Train set: papers with $\yd{p} \leq 2022$ ($n = 384$).
Test set: papers with $\yd{p} \geq 2023$ ($n = 815$).
This is the \emph{deployment-honest} protocol:
the model is trained on past literature and evaluated
on genuinely future papers.
\begin{equation}
    \operatorname{Acc}_{\mathrm{TS}}
    \;=\;
    \frac{1}{|S_{\mathrm{test}}|}
    \sum_{p \,\in\, S_{\mathrm{test}}}
    \I{\hat{d}(p) = d^*(p)}.
    \label{eq:ts}
\end{equation}
Result: $\operatorname{Acc}_{\mathrm{TS}} = 0.9804$,
macro-$F_1 = 0.9607$.

\paragraph{Protocol 2 --- Stratified $K$-Fold
Cross-Validation.}
Stratified 5-fold CV preserves per-domain proportions
across folds.
Per-fold accuracy $\operatorname{Acc}_k$ and macro-$F_1^k$
are aggregated as:
\begin{align}
    \overline{\operatorname{Acc}}_{\mathrm{CV}}
    &\;=\;
    \frac{1}{K}\sum_{k=1}^{K}\operatorname{Acc}_k,
    \label{eq:cv_acc} \\
    F_1^{\mathrm{macro}}
    &\;=\;
    \frac{1}{|\D|}
    \sum_{d \in \D}
    \frac{2\,P_d\,R_d}{P_d + R_d}.
    \label{eq:macro_f1}
\end{align}
Result: $\overline{\operatorname{Acc}}_{\mathrm{CV}}
= 0.9946 \pm 0.0040$,
macro-$F_1 = 0.9891 \pm 0.0131$.

\paragraph{Protocol 3 --- Leave-One-Source-Out (LOSO).}
For each source $s \in \mathcal{S}$, train on
$\mathcal{S} \setminus \{s\}$, test on $s$.
Mean across sources:
\begin{equation}
    \overline{\operatorname{Acc}}_{\mathrm{LOSO}}
    \;=\;
    \frac{1}{|\mathcal{S}|}
    \sum_{s \,\in\, \mathcal{S}}
    \operatorname{Acc}_s.
    \label{eq:loso}
\end{equation}
Results per source:
arXiv $0.9707$, DBLP $0.9713$,
Elsevier $1.000$, PubMed $0.9916$;
mean $\overline{\operatorname{Acc}}_{\mathrm{LOSO}} = 0.9834$,
mean macro-$F_1 = 0.9733$.

% ============================================================
\section*{Acknowledgements}
Anonymised for review.
% ============================================================

\bibliographystyle{aaai}
\begin{thebibliography}{99}

\bibitem{beel2016}
Beel, J.; Gipp, B.; Langer, S.; and Breitinger, C.
2016.
Paper recommender systems: A literature survey.
\emph{Int.\ J.\ Digit.\ Libr.}\ 17(4):305--338.

\bibitem{cohan2020}
Cohan, A.; Feldman, S.; Beltagy, I.; Downey, D.; and Weld, D.~S.
2020.
SPECTER: Document-level representation learning using
citation-informed transformers.
In \emph{ACL}, 2270--2282.

\bibitem{kleinberg1999}
Kleinberg, J.~M.
1999.
Authoritative sources in a hyperlinked environment.
\emph{J.\ ACM}\ 46(5):604--632.

\bibitem{page1999}
Page, L.; Brin, S.; Motwani, R.; and Winograd, T.
1999.
The PageRank citation ranking: Bringing order to the web.
Stanford Tech.\ Rep.\ 1999-66.

\bibitem{platt1999}
Platt, J.~C.
1999.
Probabilistic outputs for support vector machines.
In \emph{Advances in Large Margin Classifiers},
61--74. MIT Press.

\bibitem{price1965}
Price, D.~J.\ de~S.
1965.
Networks of scientific papers.
\emph{Science}\ 149(3683):510--515.

\bibitem{yao2023}
Yao, S.; Zhao, J.; Yu, D.; Du, N.; Shafran, I.;
Narasimhan, K.; and Cao, Y.
2023.
ReAct: Synergizing reasoning and acting in language models.
In \emph{ICLR}.

\bibitem{zhou2004}
Zhou, D.; Bousquet, O.; Lal, T.~N.; Weston, J.;
and Sch\"{o}lkopf, B.
2004.
Learning with local and global consistency.
In \emph{NeurIPS}, 321--328.

\end{thebibliography}

\end{document}
